\documentclass[10pt,a4paper]{amsart}
\usepackage[margin=19mm]{geometry}
\usepackage{amsmath,amssymb,mathtools}
\usepackage[scr=rsfs,cal=euler]{mathalfa}
\usepackage{xfrac}
\usepackage[unicode,hidelinks,bookmarks=false]{hyperref}
\newcommand{\ie}{i.e.\ }
\newcommand{\cf}{cf.\ }
\newcommand{\resp}{respectively\ }
\newcommand{\Subsection}{\subsection}
\providecommand{\nobreakdash}{\nobreak\hbox{-}}
\setlength{\emergencystretch}{2em}
\multlinegap0pt
\usepackage{amssymb}
\usepackage{mathtools}
\usepackage{float}
\usepackage[shortlabels]{enumitem}
\newtheorem{proposition}{Proposition}
\newcommand{\diag}{\operatorname{diag}}
\title{Algebraic Spectral Curves of Two-Channel Operator Pencils with Power-Law Response}
\author{Kejun Liu}
\date{Source: arXiv:2605.17645v4 (CC BY 4.0)}
\begin{document}
\maketitle

\noindent\emph{Source and licence.} Algebraic Spectral Curves of Two-Channel Operator Pencils with Power-Law Response. Version arXiv:2605.17645v4 (CC BY 4.0), distributed by arXiv under the Creative Commons Attribution 4.0 International licence, http://creativecommons.org/licenses/by/4.0/, as recorded in the arXiv licence metadata for this exact version and shown on the abstract page https://arxiv.org/abs/2605.17645v4. Full licence notice: \texttt{licenses/arxiv-2605.17645-CC-BY-4.0.txt}.\par\smallskip

\noindent\textit{Editorial scope.} Counted unit: the article's proposition prop:higher-channel-EJ (source0.tex lines 1867--1880). The article's complete original proof of it is delivered with it (source0.tex lines 1882--1924, one \texttt{proof} environment of 43 lines). No mathematics is written, repaired or re-derived: the statement and the proof are the article's own lines, in the article's own notation, and no retained line is substituted or re-flowed.  Retained uncounted as the source-local context the proof needs: lines 1844--1846.  Nothing was omitted from any retained span, so the package carries no omission record.  No retained line contains a citation, so the wrapper carries no bibliography and no bibliography entry.  No retained line contains a figure environment, an image-inclusion command or drawing code, so no image is delivered and the package declares no asset dependency.  Editorial formatting only, in addition: an AMS \texttt{amsart} preamble that restates the article's own declarations and macros used by the retained lines, namely the environments \texttt{proposition} on separate counters, together with the standard packages those lines need.  The retained environments print in this document as Proposition 1 in order of appearance, whereas the article prints the same items with the numbering of its own full document.  The complete native source of the article as distributed in the arXiv e-print for this exact version is bundled unchanged as \texttt{sources/200807/source0.tex}.
\begin{equation}\label{eq:higher-channel-curve}
 \det(u^nI_N-H_0-wV)=0.
\end{equation}

\begin{proposition}[Indefinite two-energy slice]\label{prop:higher-channel-EJ}
Let
\[
 J=\diag(I_{p_+},-I_{p_-}),\qquad
 p_+,p_-\geq1,\qquad p_++p_-=N,
\]
and set $H_0=EJ$ with $E\ne0$.  For every $n\geq2$, there is a
nonempty Zariski-open subset of matrices $V=JS$, $S^\top=S$, on which
the normalization of \eqref{eq:higher-channel-curve} is geometrically
integral and has genus
\begin{equation}\label{eq:higher-channel-EJ-genus}
 \boxed{g_{EJ}=np_+p_--N+1.}
\end{equation}
\end{proposition}

\begin{proof}
First consider the degree-$N$ projective base curve
\[
 B:\quad \det(XI_N-EZJ-WV)=0.
\]
On $W=0$ it has the two forced points
$P_+=(E:0:1)$ and $P_-=(-E:0:1)$.  Relative to the positive and negative
blocks of $J$, their tangent cones are, up to nonzero factors,
\[
 \det(\xi I_{p_+}-WV_{++}),\qquad
 \det(\eta I_{p_-}-WV_{--}).
\]
For generic $V$, the two diagonal blocks are regular semisimple.  Hence
$P_+$ and $P_-$ are ordinary multiple points of multiplicities $p_+$ and
$p_-$.  The required generic locus is nonempty: start with a diagonal
$V$, whose same-sign eigenvalue lines meet only at the two forced points,
and perturb its symmetric representative $S=JV$ by generic off-diagonal
entries.  This smooths the cross-sign intersections while preserving the
two ordinary multiple points.  The cross-sign incidence graph is connected,
so smoothing all its edges makes the nearby fiber irreducible.  No further
singularity persists on a generic projective fiber.

The normalization $B^\nu$ consequently has genus
\begin{align*}
 g(B^\nu)
 &=\frac{(N-1)(N-2)}2
   -\frac{p_+(p_+-1)+p_-(p_--1)}2 \\
 &=(p_+-1)(p_--1).
\end{align*}
Equivalently, this is Vivolo's normal-crossing correction after the
projective change $t=1/W$, $y=X/W$, which writes the base curve as
$\det(V+tEJ-yI_N)=0$.

On the same generic locus, $x=X/Z$ has $N$ simple zeros and $N$ simple
poles on $B^\nu$.  Equation \eqref{eq:higher-channel-curve} is the
cyclic pullback $u^n=x$, totally ramified at those $2N$ points and nowhere
else.  Riemann--Hurwitz gives
\[
 2g_{EJ}-2
 =n\bigl(2(p_+-1)(p_--1)-2\bigr)+2N(n-1),
\]
which reduces to \eqref{eq:higher-channel-EJ-genus}.
\end{proof}

\end{document}
